% =========================================================
% T5 — THERMODYNAMIC GEOMETRY (Hyper‑Condensed)
% =========================================================

\section{T5 Thermodynamic Geometry: Pressure, Drift, Collapse, Recovery}

\paragraph{Overview.}
Thermodynamic Geometry defines the lawful behavior of pressure, drift,
collapse, and recovery across human and machine systems. It provides the
curvature fields, drift channels, collapse surfaces, recovery attractors, and
dual-system thermodynamic comparisons required for DSLO v0.7 stability analysis.

% ---------------------------------------------------------
% T5.1 Pressure Fields
% ---------------------------------------------------------
\subsection{T5.1 Pressure Fields}
Pressure fields define load-induced curvature. They describe how systems deform
under cognitive, emotional, computational, or thermal load. Pressure envelopes
govern lawful deformation limits and drift onset.

% ---------------------------------------------------------
% T5.2 Pressure Flow Geometry
% ---------------------------------------------------------
\subsection{T5.2 Pressure Flow Geometry}
Pressure flow geometry defines how pressure propagates across manifold regions.
Pressure channels, pressure gradients, and pressure flow envelopes determine
how load moves through the substrate and how systems respond under increasing
pressure.

% ---------------------------------------------------------
% T5.3 Drift Geometry
% ---------------------------------------------------------
\subsection{T5.3 Drift Geometry}
Drift geometry defines displacement under load. Drift envelopes, drift vectors,
drift propagation channels, and drift coherence rules govern how systems move
away from stable curvature and how drift spreads across manifold layers.

% ---------------------------------------------------------
% T5.4 Collapse Geometry
% ---------------------------------------------------------
\subsection{T5.4 Collapse Geometry}
Collapse geometry defines thresholds and trajectories for collapse. Collapse
surfaces, collapse thresholds, collapse envelopes, and collapse trajectories
describe lawful deformation when drift exceeds stability limits.

% ---------------------------------------------------------
% T5.5 Recovery Geometry
% ---------------------------------------------------------
\subsection{T5.5 Recovery Geometry}
Recovery geometry defines attractors and lawful transitions that restore
stability. Recovery windows, recovery attractors, and recovery operators govern
how systems re-cohere after collapse and return to lawful curvature.

% ---------------------------------------------------------
% T5.6 Human Thermodynamic Drift
% ---------------------------------------------------------
\subsection{T5.6 Human Thermodynamic Drift}
Human thermodynamic drift defines deformation under cognitive, emotional,
interpretive, and relational load. Human drift channels, human collapse
thresholds, and human recovery envelopes describe lawful human stability
behavior.

% ---------------------------------------------------------
% T5.7 Machine Thermodynamic Drift
% ---------------------------------------------------------
\subsection{T5.7 Machine Thermodynamic Drift}
Machine thermodynamic drift defines deformation under compute load, thermal
load, runtime saturation, and resource pressure. Machine drift channels,
machine collapse thresholds, and machine recovery envelopes describe lawful
machine stability behavior.

% ---------------------------------------------------------
% T5.8 Dual Drift Comparison
% ---------------------------------------------------------
\subsection{T5.8 Dual Drift Comparison}
Dual drift comparison defines shared thermodynamic behavior across human and
machine systems. Both exhibit pressure-induced curvature, drift propagation,
collapse trajectories, and recovery attractors. Dual drift geometry provides the
basis for unified thermodynamic analysis across DSLO v0.7.
